\chapter{多目标与受约束决策}
\label{chap:rs-constraints}

同一个推荐页面可以得到不同的评价。用户在通勤背包中找到了合适的款式，商家获得了一次展示，平台收到了广告费；如果页面反复出现相似广告，用户也可能更快离开。前面几章已经给出了候选、响应预测和展示方案，现在还缺少一个判断：哪些结果值得追求，哪些代价可以接受，哪些选择根本不允许执行？一个预测准确的模型不会自动回答这些问题。

本章把目标和约束放在同一个决策问题中。目标说明希望获得什么，约束说明允许作出什么选择；两者都必须明确受益主体、计量单位和时间范围。在这个共同定义上，加权、优先级和目标约束表达不同的取舍原则。累计注意、双边分配和集合可靠性进一步说明，看似相同的“兼顾多方”可以对应很不一样的数学要求。得到方案后，还要分别检查可行性、目标改善与分群后果，才能确定应调整问题定义、预测模型还是求解过程。

\section{决策问题定义}
\label{sec:rs-decision-definition}

沿用第~\ref{chap:rs-presentation}~章的记号，$a_t$表示时点$t$的一次展示行动，可以包含对象列表、位置、表达、发送时间和交互方式。给定请求信息$\bm{x}_t$，候选目录与资格规则确定行动范围$\mathcal{A}_t$；库存、预算和服务能力又把它收缩为可行范围$\mathcal{F}_t$。对同一个行动定义多个目标
\[
 \bm{f}(a_t\mid\bm{x}_t)
 =\bigl(f_1(a_t\mid\bm{x}_t),\ldots,
        f_d(a_t\mid\bm{x}_t)\bigr).
\]
这里$d$为目标数，各分量暂统一为越大越好；重复、投诉等希望降低的量，可以取负号，也可以保留为上界约束。后续求解必须同时使用$\bm{f}$和$\mathcal{F}_t$。先算一个最大分数、上线前再随意删改，通常已经改变原来求解的问题。

目标有时作用于一次请求，有时作用于一批用户或一段时间。若要求某位提供者本周至少获得一定机会，就不能只把$f_d$写成当前用户的响应。决策单位可以扩大为一批分配，也可以保留逐请求执行、显式维护累计状态。应由要求的作用范围决定变量，不由现成模型输出的形状决定问题。

\subsection{多方目标}
\label{sec:rs-multistakeholder-objectives}

用户关心能否完成任务，平台关心服务收入与持续使用，内容提供者关心触达和消费，广告主关心投放带来的价值。这些主体可能分享某些目标，也可能在同一页面上发生冲突。例如，展示更昂贵的背包可能提高预计成交金额，却减少满足500元预算的选项；给新商家更多机会，也可能减少当前高相关商家的展示。

将这些诉求变成目标，需要说明观察什么。表~\ref{tab:rs-objective-contracts}列出可用于讨论的测量口径。表中的行为是证据来源，不能凭名称断言它们已经完整测量了相应利益。一次长停留可能来自认真比较，也可能来自信息难找；点击广告后购买，也不说明没有广告就不会购买。第~\ref{sec:rs-uplift-ranking}~节建立的增量区别在这里仍然适用。

\begin{table*}[htbp]
 \centering
 \small
 \begin{tabularx}{\linewidth}{@{}p{19mm}p{31mm}X X@{}}
  \toprule
  关注主体 & 可观察的量 & 必须固定的口径 & 不能直接替代的结果 \\
  \midrule
  用户 & 任务完成、问卷满意度 & 任务定义、纳入用户、观察窗口 & 点击、停留不自动等于满意 \\
  平台 & 收入、服务成本 & 请求或用户口径，毛收入或净收益 & 成交金额不等于平台收入 \\
  提供者 & 可见曝光、有效消费 & 物品或创作者，位置注意与累计周期 & 被选入候选不等于被看见 \\
  广告主 & 转化价值、实际支出 & 转化成熟窗口、价值单位、费用口径 & 归因转化不等于增量转化 \\
  \bottomrule
 \end{tabularx}
 \caption{多方目标的测量约定。每行只是可能的指标入口；定义目标时，还须说明估计依据、受益对象以及允许牺牲的幅度。}
 \label{tab:rs-objective-contracts}
\end{table*}

一次推荐的响应预测可以为决策提供输入，但预测量与决策价值有不同含义。设$p_{\mathrm{click}}(a)$是页面发生点击的预测概率，$v(a)$是点击后带来的价值估计，$c(a)$是执行成本；只有在这些条件概率与价值定义相容时，才能用$p_{\mathrm{click}}(a)v(a)-c(a)$表示相应的预期净价值。若$v(a)$本身已经是按展示平均的价值，再乘一次点击概率就会重复折算。多任务模型有多个输出头，也不表示每个输出应该同权进入最终分数。

短期目标可写成当前反馈的条件期望$f_\ell(a_t\mid\bm{x}_t)$。累计目标则需要给出周期和聚合方式，例如$T$次请求中的总收入$\mathbb{E}[\sum_{t=1}^{T}r_t]$，或每位用户在一个月内完成任务的比例。平均每次请求、平均每位用户和总量会对高频用户赋予不同权重。若行动会改变未来请求、用户状态或供给，累计目标还依赖后续策略；第~\ref{chap:rs-long-term}~章将把这种依赖写入连续决策。

提供者公平也要选定比较标准。按物品平均曝光、按创作者平均曝光、按相关性比例分配注意，分别对不同对象作出承诺。一位创作者有十个物品，另一位只有一个，逐物品相同曝光并不等于逐创作者相同曝光。相关性比例标准又依赖相关性估计是否可信；若新提供者缺少历史，低预测值可能同时反映信息不足。目标定义应交代这种证据边界，不能让一个未解释的“公平分数”同时承担所有标准。

还须在求解前写清各方允许交换什么。提高收入是否可以降低任务完成率，若可以，最多降低多少？提高总体覆盖是否允许某个小群体的体验持续下降？这些是决策要求，不会由点击率和金额的数值大小自然确定。下一节的约束可以把不能交换的要求从目标中单独表示出来。

\subsection{资源与约束}
\label{sec:rs-decision-resources}

可行行动必须符合资格和资源条件。对象已下架、用户不可访问、同一列表不能重复等规则，决定基本可选范围；库存和服务容量限制能交付多少；广告预算限制周期内能够支出多少。若限制写成$g_j(a_t,\bm{b}_t)\leq0$，其中$\bm{b}_t$保存请求开始前的资源状态，就得到
\begin{equation}
 \mathcal{F}_t=
 \{a\in\mathcal{A}_t:
   g_j(a,\bm{b}_t)\leq0,\ j=1,\ldots,m\}.
 \label{eq:rs-feasible-actions}
\end{equation}
$m$为约束数。约束所用信息必须在决策时可用；事后才知道的消费或成本，不能当成当时已经确定的数值。

\term{硬约束}（hard constraint）直接排除违反条件的方案。\term{软惩罚}（soft penalty）则允许违反，只是在目标中扣除相应代价。例如“同一商家最多两条”可写成硬计数上限；“超过两条后每增加一条扣除$\lambda$分”则是软惩罚。后一种规则可能因其他收益足够高而选择第三条。阈值这个词本身不能区分两者，必须看求解器是否允许越过阈值。

公平要求同样可以进入可行范围。若要求一组提供者在整个周期的可见曝光不低于$E_{\min}$，应约束该组的累计曝光；若要求两组平均机会的差距不超过$\epsilon$，应说明平均的分母和允许资格。把差距写成要最小化的目标，则表示愿意与其他目标交换，没有自动给出$\epsilon$保证。这两种表述可以同时存在，例如先限制最大差距，再在其内继续改善公平。

跨请求资源会随执行改变。设$b_t$为请求开始时的预算余额，$c_t$为本次最终计费的实际成本，在没有退款和充值的简化情形下，
\[
 b_{t+1}=b_t-c_t.
\]
预测成本$\widehat c_t$可以用于计划，但不能代替账本中的$c_t$。成本只有在点击或转化后才确定时，决策还需要预留额度、并发计费控制或其他明确的执行机制。仅约束期望支出不超过预算，不能推出每条实现路径都不超支；广告中的具体控制见第~\ref{chap:rs-advertising}~章。

\begin{example}[当前可行与周期可行]
 \label{ex:rs-two-request-budget}
 假设一个已知只含两次展示的周期共有10元预算，每次展示有两个对象位置，费用在执行时确定。还要求商家甲在周期内至少占两个位置，且同一页面至多出现一次。

 第一次若选择6元、包含甲一次的方案，执行后余额为4元，甲已获得一次曝光。第二次包含甲的最低费用是5元，于是第一次虽未超预算，余下要求已经无法同时完成。若第一次改选4元、同样包含甲一次的方案，第二次再花5元展示甲，周期结束时支出9元、余额1元，甲累计获得两次曝光。
\end{example}

这个例子使用了已知的第二次机会。实际未来流量和可选对象通常不确定，此时至少应区分“按预测仍有完成空间”与“对所有允许的未来变化都可完成”。对最低机会约束，还可检查剩余缺口是否超过未来可分配位置数；对预算，则要检查为最低承诺预留的成本。简单的剩余容量检查未必充分，但能及早发现只看当前分数掩盖的冲突。

没有可行行动时，求解器应报告冲突的资格、容量或阈值，并按预先定义的降级规则处理。若允许缩短列表、延后请求或放宽某项软要求，应把这种行动写进问题。直接返回最高分的不可行列表，会把目标冲突误装成模型选择。

\section{决策方案求解}
\label{sec:rs-decision-solvers}

目标和可行范围确定后，仍可能有多个互不胜出的方案。下面固定三个假设页面，参与度与收入分别为A$(0.9,1)$、B$(0.8,3)$、C$(0.6,5)$。参与度表示一次请求内发生有效参与的概率，收入单位为元／请求；数值只用于教学，并假设三个方案均满足共有资源约束。收入增加伴随参与度降低，因此两项信息不足以选出唯一方案。

加权方法规定一次目标损失可以由多少另一目标补偿；优先级方法规定较低优先目标何时有资格改变选择；目标约束方法则规定必须达到什么底线。它们可以组合使用。例如，先满足预算和参与度底线，再最大化收入与内容差异的加权和。表~\ref{tab:rs-objective-coordination}给出贯穿例子的直接结果，后文展开每一项计算。

\begin{table}[htbp]
 \centering
 \small
 \begin{tabularx}{\linewidth}{@{}p{24mm}X p{14mm}@{}}
  \toprule
  原则 & 本例采用的规则 & 选择 \\
  \midrule
  权重 & 最大化参与度$+0.08\times$收入 & B \\
  严格优先级 & 参与度优先，收入只打破参与度平局 & A \\
  容忍优先级 & 参与度距最佳值不超过0.1，再提高收入 & B \\
  目标约束 & 参与度至少0.8，再最大化收入 & B \\
  \bottomrule
 \end{tabularx}
 \caption{在同一可行方案集合上比较目标协调原则。相同选择不表示规则相同；候选范围或目标值改变后，各规则可能分离。}
 \label{tab:rs-objective-coordination}
\end{table}

直接求解的输出是一个行动或一批分配。若训练一个模型来生成方案，输出则先是模型参数，随后还要经过评分、列表构造和执行检查。训练问题满足条件，不自动表示后续每个行动满足条件，尤其不能忽略截断、去重和资源更新对结果的改变。

\subsection{基于权重的多目标优化}
\label{sec:rs-weighted-objectives}

在目标允许相互补偿时，可以选定非负权重$w_\ell$，最大化加权目标
\begin{equation}
 a^\star\in\argmax_{a\in\mathcal{F}_t}
     \sum_{\ell=1}^{d}w_\ell f_\ell(a).
 \label{eq:rs-weighted-decision}
\end{equation}
权重同时处理单位转换和偏好强弱。若用$f_{\mathrm{eng}}+\lambda f_{\mathrm{rev}}$，$\lambda$就表示一元收入在当前比较中相当于多少参与度。它不是两个模型都输出0到1时自然产生的常数。

对A、B、C，$\lambda=0.02$时，分数依次为$0.92,0.86,0.70$，选择A；$\lambda=0.1$时为$1.0,1.1,1.1$，B与C同分；若预先规定同分时优先参与度较高的方案，就选择B。取$\lambda=0.2$时则为$1.1,1.4,1.6$，选择C。平局处理也是决策规则的一部分，不能在得到结果后临时指定。

比较两两分差还可确定权重的作用区间。B胜过A要求$-0.1+2\lambda>0$，即$\lambda>0.05$；B胜过C要求$0.2-2\lambda>0$，即$\lambda<0.1$。所以$0.05<\lambda<0.1$时B是唯一的最优选择。在边界上应按既定平局规则处理。把收入从元换成分后，数值变为$100,300,500$，相同偏好要求把$\lambda$除以100；若仍用原权重，改变的是取舍。

标准化也不会消除这一步选择。按训练集标准差缩放，或按历史最好值归一化，都会引入参考群体和时间范围。参考分布改变后，同一权重可能对应不同的实际交换率。实际选权重可以从可接受代价、候选权重下的验证结果及多方要求出发；应同时保留原始单位的指标，让决策者知道组合分数提高由什么贡献。

将软惩罚加入式~\eqref{eq:rs-weighted-decision}很方便，例如对超过来源份额的量扣分。然而，对任意有限惩罚系数，若其他收益没有足够小的已知上界，就可能出现值得违反的方案。需要绝不越过的库存、资格和费用条件仍应留在$\mathcal{F}_t$中。拉格朗日乘子可以帮助求解约束问题，其对偶最优性和恢复可行解的条件，见\BookChapterRef{chap:optimization}{最优化}。

加权训练与加权决策也要区分。训练损失的权重$\omega_\ell$决定共享参数怎样拟合不同任务，决策权重$w_\ell$决定根据预测如何选行动。加大点击损失的权重，可能提高点击预测精度，但它并不等于在上线目标中提高点击的交换率。第~\ref{chap:rs-ranking}~章的MMoE、PLE和PEPNet主要组织响应估计；UniR$^2$给出多个目标的分数，PinRec通过候选名额改变覆盖。它们都需要与明确的决策目标和可行范围衔接。

\subsubsection{PE-LTR：偏好约束的Pareto排序学习}
\label{sec:rs-pe-ltr}

固定一组训练权重时，不同任务的梯度可能方向相反。沿点击损失下降，价值损失可能上升；权重不变也未必适合整个训练过程。\term{PE-LTR}（Pareto-efficient learning to rank）根据各目标当前的梯度调整组合权重，并通过权重下界限制允许的偏好范围。它学习一个排序模型，再比较不同偏好设定训练出的方案\scite{rs-r90}

设共享参数为$\bm{\theta}\in\mathbb{R}^{p}$，有$d$个可微损失$\mathcal{L}_\ell(\bm{\theta})$，梯度为$\bm{g}_\ell=\nabla_{\bm{\theta}}\mathcal{L}_\ell$。论文将权重选择建立为最小范数组合问题：
\begin{equation}
 \begin{aligned}
  \min_{\bm{\omega}\in\mathbb{R}^{d}}\quad&
   \left\lVert\sum_{\ell=1}^{d}\omega_\ell\bm{g}_\ell
   \right\rVert_2^2,\\
  \text{s.t.}\quad&
   \sum_{\ell=1}^{d}\omega_\ell=1,\qquad
   \omega_\ell\geq c_\ell.
 \end{aligned}
 \label{eq:rs-pe-gradient}
\end{equation}
下界$c_\ell\geq0$必须满足$\sum_\ell c_\ell\leq1$，否则没有合法权重。这里的归一化作用于权重和，并不自动消除各损失及其梯度的尺度差异；改变金额单位或损失缩放仍会改变解。

给定权重，用加权梯度更新
\[
 \bm{\theta}^{(s+1)}
 =\bm{\theta}^{(s)}
  -\eta\sum_{\ell=1}^{d}\omega_\ell^{(s)}
       \nabla\mathcal{L}_\ell(\bm{\theta}^{(s)}),
\]
再用新参数下的梯度计算后续权重。$\eta>0$为学习率，$s$为训练迭代。原文的求解器先保留等式约束，求一个放松解，再通过投影取得合法权重。理解这一实现时要保留原问题的度量：若$\bm{G}$以各梯度为行，目标是$\bm{\omega}^{\top}\bm{G}\bm{G}^{\top}\bm{\omega}$，一般不等于权重空间中的欧氏距离。投影后满足上下界，只能直接说明权重可行，不能仅凭这一步断言已经精确求解原二次规划。小目标数时可直接求解式~\eqref{eq:rs-pe-gradient}，并检查约束和最优性残差。

两目标的情形可以手算。令$\omega_1=\alpha$、$\omega_2=1-\alpha$，当$\bm{g}_1\ne\bm{g}_2$时，精确解为
\[
 \alpha^\star=
 \operatorname{clip}_{[c_1,\,1-c_2]}
 \left(
 \frac{(\bm{g}_2-\bm{g}_1)^\top\bm{g}_2}
      {\lVert\bm{g}_1-\bm{g}_2\rVert_2^2}
 \right).
\]
若梯度相同，任意合法权重都产生相同方向。假设点击与价值损失梯度分别为$(1,0)$和$(-1,1)$，不设正下界时$\alpha^\star=0.6$，组合梯度为$(0.2,0.4)$。它与两个原梯度的内积都为0.2，所以沿其负方向走足够小一步，两项损失在一阶近似下都降低。若要求点击权重至少0.8，组合变为$(0.6,0.2)$；此时它与价值梯度的内积为$-0.4$，沿负方向更新反而会提高价值损失。这正是偏好限制可能付出的代价。

这个差别也限定理论解释。当$c_\ell=0$时，精确的最小范数组合若非零，可给出共同下降方向；若为零，则达到不存在共同严格下降方向的一阶条件，见\BookSectionRef{sec:opt-multiobjective-saddle}{多目标与鞍点}。增加权重下界后，允许的梯度组合范围变小，非零解不再普遍保证所有损失同时下降。即使组合为零，一般非凸网络中也只能作一阶驻点判断，不能据此声称全局Pareto最优。随机小批量梯度和有限训练进一步要求在独立数据上检查实际结果。

PE-LTR的电商实例使用曝光、点击与购买信息定义点击和价值相关损失。点击项采用点击标签的对数损失；价值项以价格的单调变换加权购买概率的对数。原文把购买分解为点击与点击后购买，并在该实例中假定点击后购买概率不依赖当前排序参数，使价值梯度主要落到价格加权的点击概率上。这是特定代理目标，并不是直接对真实GMV求导；若排序改变了点击后的购买人群，这种简化需要重新检查。框架也可使用其他可微损失，但损失来源和业务意义必须重新说明。

训练后保留排序参数，用它对请求候选评分并生成列表。改变$c_\ell$、重复训练，可以获得一组不同取舍的候选模型；比较时应报告同一数据、候选范围和资源条件下的点击与价值，而不只展示训练损失。计算代价包括各任务的梯度，以及目标数为$d$的权重子问题；直接形成梯度Gram矩阵需要约$O(d^2p)$计算和各梯度的存储。任务数小不表示反向传播免费。

最重要的接口边界是：$\omega_{\mathrm{click}}\geq0.8$约束的是训练权重，它不表示点击率至少0.8，也不表示点击率相对基线损失不超过某个幅度。若要作后一类承诺，就要进入目标阈值、统计证据和实际执行的检查。

\subsection{基于优先级的多目标优化}
\label{sec:rs-priority-objectives}

有些目标不接受任意补偿。例如，在资格合规且资源足够的页面中，希望尽量保留任务完成率，收入只在相同质量的方案之间起作用。\term{严格词典序}（lexicographic order）表达这种规则：按既定次序比较目标，在第一个不同的目标上决定优劣，后面的改进不能抵消前面的下降。

令$\mathcal{F}_0=\mathcal{F}_t$，$f_1,\ldots,f_d$已经按优先级排列。逐层求
\[
 f_\ell^\star=\max_{a\in\mathcal{F}_{\ell-1}}f_\ell(a),
 \qquad
 \mathcal{F}_\ell=
 \{a\in\mathcal{F}_{\ell-1}:f_\ell(a)=f_\ell^\star\}.
\]
最终从$\mathcal{F}_d$选择行动；完全同分时仍要规定可复现的处理方式。对A、B、C，参与度优先会在第一层只留下A，收入不再改变结果。若收入优先，则第一层选择C。这是两种不同要求，不能只写“按照优先级排序”而不说明次序。

严格相等在预测有误差或连续目标下可能过于刚性。允许高优先级目标距最佳值损失不超过$\epsilon_\ell$时，可以改为
\[
 \mathcal{F}_\ell=
 \{a\in\mathcal{F}_{\ell-1}:
       f_\ell(a)\geq f_\ell^\star-\epsilon_\ell\}.
\]
本例令参与度容忍量为0.1，最佳值0.9给出门槛0.8，留下A、B，再按收入选择B。这里0.1是概率值的绝对差，即10个百分点，不是相对降低10\%。容忍量应随目标口径一起记录。

有限大权重有时能模仿词典序，但需要额外条件。若高优先目标的不同值至少相差$\Delta>0$，而低优先目标的最大差为$M$，让前者权重足够大于$M/\Delta$即可压过任何低优先改善。连续目标可能没有正的最小差，范围未知时也无法确定足够大的权重。直接逐层求解则保留了规则本身；近似求解时，还应把各层最优值误差纳入容忍量解释。

优先级可以作用于目标组，例如体验组优先于收入组。此时还须规定组内比较方法，否则只把冲突移到了组内。资源约束仍然在每一层生效，参与度再高也不能突破预算；若优先级来自业务治理要求，应当明确记录，而不是从某个训练损失的大小倒推。

\subsection{基于目标约束的多目标优化}
\label{sec:rs-goal-constrained-objectives}

另一种常见要求是给某些目标规定绝对底线，在其余空间内改善主目标。例如
\begin{equation}
 \begin{aligned}
  \max_{a\in\mathcal{F}_t}\quad&f_{\mathrm{rev}}(a),\\
  \text{s.t.}\quad&f_{\mathrm{eng}}(a)\geq\tau
 \end{aligned}
 \label{eq:rs-target-constraint}
\end{equation}
将参与度的可接受水平写成$\tau$。对三个方案，$\tau=0.8$时A、B可行，选择B；$\tau=0.85$时仅A可行；$\tau=0.6$时选择C；$\tau>0.9$时没有可行方案。阈值改变可行集合，求解器不能用一个很高的收入分数补偿资格缺失。

本例中，参与度允许距最佳值损失0.1，恰好与绝对底线0.8得到相同结果。但候选集合改变后两者会分离：若最佳参与度降为0.85，前者的门槛随之变为0.75，后者仍为0.8。相对最佳方案的容忍、相对历史基线的变化和绝对阈值，都有自己的参照物。

目标约束也可以作用于一批分配或总体风险。前者可以直接检查某个确定分配是否满足数量条件；后者涉及未知数据分布，需要校准或实验提供统计证据。若用预测参与度代替真实参与度，式~\eqref{eq:rs-target-constraint}最多直接保证预测量达标。即便求解精确，预测误差仍可能让实际指标越界。下面三个实例分别处理累计注意、双边集合分配和集合错误风险，条件与保证不能互换。

\subsubsection{Equity of Attention：累计注意分配}
\label{sec:rs-equity-attention}

两个同样相关的对象不能同时占据唯一首位，若位置注意不同，单次排列就未必能按相关性比例分配机会。\term{Equity of Attention}把比较窗口延长到多次排序：记录每个对象累计获得的注意和累计相关性，在后续请求中补偿过去的差距，同时限制本次排序质量的损失\scite{rs-r17}

设本次有$n$个待排序对象，归一化相关性为$r_{i,t}\geq0$，满足$\sum_i r_{i,t}=1$；位置$j$的注意份额为$w_j\geq0$，满足$\sum_j w_j=1$。这两个量分别描述应获得的比例与位置实际分配的比例。若页面只关注前$k$位，可把更后面的位置注意设为零，但仍需说明截断与归一化。相关性是评价相关程度的估计，$w_j$则需要位置查看等证据；未经纠偏的点击率不能同时无条件充当两者。

令$\bm{X}_{i,j}\in\{0,1\}$表示把对象$i$放到位置$j$，并要求每个对象和位置恰好匹配一次。请求前的累计注意为$A_{i,t-1}$，累计相关性为$R_{i,t-1}$，执行后更新为
\[
 A_{i,t}=A_{i,t-1}+\sum_jw_j\bm{X}_{i,j},
 \qquad
 R_{i,t}=R_{i,t-1}+r_{i,t}.
\]
初始值可取零；不在本次比较范围的对象应按已声明的资格口径处理，不能把不可供给时间默认为获得了新的相关性份额。原方法用$\sum_i|A_{i,t}-R_{i,t}|$衡量累计差距。

为限制当前质量，使用相对于相关性降序排列的归一化折损累计增益（normalized discounted cumulative gain, NDCG）。这里先把它当作质量约束的具体计算：位置$j$的相关性收益为$(2^{r_{i,t}}-1)/\log_2(j+1)$，理想排列的前$k$位总收益记作$\operatorname{IDCG}_t@k$。只在该值大于零时使用比值形式。给定阈值$\tau\in[0,1]$，整数指派问题可写成
\begin{equation}
 \begin{aligned}
  \min_{\bm{X}}\quad&
    \sum_{i=1}^{n}\sum_{j=1}^{n}
       \left|A_{i,t-1}+w_j-R_{i,t-1}-r_{i,t}\right|
       \bm{X}_{i,j},\\
  \text{s.t.}\quad&
    \sum_{j=1}^{k}\sum_{i=1}^{n}
       \frac{2^{r_{i,t}}-1}{\log_2(j+1)}\bm{X}_{i,j}
       \geq\tau\,\operatorname{IDCG}_t@k,\\
  &\sum_i\bm{X}_{i,j}=1,\quad
   \sum_j\bm{X}_{i,j}=1,\quad
   \bm{X}_{i,j}\in\{0,1\}.
 \end{aligned}
 \label{eq:rs-equity-assignment}
\end{equation}
绝对值内的量对每个对象—位置组合都已确定，因此目标对$\bm{X}$是线性的。求解器输出排列后再更新累计状态。NDCG的通用评价口径见第~\ref{chap:rs-evaluation}~章；这里使用当前相关性分数作质量参照，不表示已经获得真实用户满意度的保证。

假设两个对象A、B每次相关性均为$(0.5,0.5)$，位置注意为$(0.8,0.2)$。第一轮把A放前面后，累计差距是$|0.8-0.5|+|0.2-0.5|=0.6$。第二轮若顺序不变，累计注意为$(1.6,0.4)$、累计相关性为$(1,1)$，差距增至1.2；若交换位置，累计注意变为$(1,1)$，差距归零。两个对象相关性相同，所以两轮的NDCG都是1，交换不损失这个质量指标。图~\ref{fig:rs-attention-accounting}显示了两种累计结果。

\begin{figure}[htbp]
 \centering
 % 教学假设：两对象每轮相关性均为0.5，位置注意为0.8/0.2。
\begin{tikzpicture}[
 x=1mm,y=1mm,font=\figurefont\small,
 box/.style={draw=BookBlue,line width=0.85pt,fill=BookBlue!5,
             minimum width=44mm,minimum height=13mm,align=center},
 flow/.style={->,>=latex,line width=1.1pt,draw=BookInk,
              rounded corners=1.2mm}
]
 \path[use as bounding box] (0,0) rectangle (112,65);
 \node[box] (first) at (56,56) {第一轮：A在前\\累计注意$(0.8,0.2)$};
 \node[box] (same) at (25,25) {第二轮：顺序不变\\累计注意$(1.6,0.4)$};
 \node[box,draw=BookTeal,fill=BookTeal!5] (swap) at (87,25)
    {第二轮：交换位置\\累计注意$(1,1)$};
 \draw[flow] (first.south) -- (56,41) -| (same.north);
 \draw[flow] (56,41) -| (swap.north);
 \node[text=BookMuted,font=\figurefont\footnotesize] at (56,38)
    {两轮累计相关性均为$(1,1)$};
 \node[text=BookInk] at (25,7) {累计差距$1.2$};
 \node[text=BookTeal] at (87,7) {累计差距$0$};
\end{tikzpicture}

 \caption{两次排序的累计注意账本。教学假设每次相关性份额为$(0.5,0.5)$、首末位置注意为$(0.8,0.2)$；交换第二次顺序可补偿第一轮的差距。注意份额是模型规定的期望分配，不是每个真实用户必然分配的观看时间。}
 \label{fig:rs-attention-accounting}
\end{figure}

若A、B相关性不同，交换可能使NDCG低于阈值，求解就要在允许的排列中寻找更小差距。$\tau=1$只允许不降低质量的安排，仍可在相关性平局内重排。没有额外资格限制且保留完整原排列时，原排列本身提供一个可行解；加入候选裁剪、来源限制后，应重新核验这一点。不能因为本轮选择了最小差距，就承诺任意未来流量下累计差距必定收敛到零。

这个问题含$n^2$个二元指派变量，质量约束使它不能简化为普通排序。原文先筛选较小的对象集合以减少在线求解规模，但筛选也可能限制后续补偿。实际应记录求解时限、可行解与最优界的差距。若页面注意受内容和邻接对象影响，固定$w_j$的模型也可能失真。

累计注意按相关性匹配是一种明确的比例标准。只有相关性相同时，它才退化为相同注意；即便逐对象都满足比例关系，也不能解释为每位创作者获得相同曝光，更不直接保证创作收入或存续。供给者对机会的响应由第~\ref{chap:rs-long-term}~章继续分析。

\subsubsection{FairRec：双边公平的推荐分配}
\label{sec:rs-fairrec}

累计注意处理对象在不同位置获得的机会。若先忽略位置，把一批用户各自得到的推荐集合共同分配，则还会出现另一种冲突：为缺少曝光的产品补机会，可能让部分用户承担更多相关性损失。\term{FairRec}把推荐写成不可分对象的分配问题，用用户轮流选择的过程协调用户效用与产品最低曝光\scite{rs-r19}

设用户集合$U$有$m$人，产品集合$I$有$n$项，每人得到$k$个不同产品。相关性矩阵$\bm{V}\in\mathbb{R}_{\geq0}^{m\times n}$在分配前固定，用户$u$对集合$S$的效用取可加形式
\[
 v_u(S)=\sum_{i\in S}\bm{V}_{u,i}.
\]
集合$S_u$是分配结果，产品$i$的曝光为$e_i=\sum_u\mathbb{I}(i\in S_u)$。这里假设集合内各对象得到相同的曝光计数，没有位置折扣；因此总曝光为$mk$。相关性可以来自前文的模型，但FairRec本身不重新训练这个矩阵。

用户侧采用移除一个对象后的无嫉妒（envy-free up to one good, EF1）。其含义是：用用户$u$自己的评价比较两个集合，即使$u$更喜欢他人的集合，从那个集合中移去某一个对象后，这种嫉妒就消失。对任意$u,w$且$S_w\ne\varnothing$，要求存在$j\in S_w$使
\begin{equation}
 v_u(S_u)\geq v_u(S_w\setminus\{j\}).
 \label{eq:rs-fairrec-ef1}
\end{equation}
等式两侧都使用$v_u$，不能拿用户$w$的评分评价被移除对象。EF1允许原集合仍被嫉妒，也没有要求每位用户得到其独立Top-$k$的固定比例。

产品侧以\term{最大最小份额}（maximin share, MMS）作为最低曝光的参照：设想把总曝光尽量均分后，每个产品可以得到的整数份额。在本问题的均匀曝光口径中，这个参照为
\[
 \ell=\left\lfloor\frac{mk}{n}\right\rfloor.
\]
原论文在$k<n\leq mk$、份额参数取1时建立保证；此时$\ell\geq1$。它不承诺所有产品都取得$\ell$份，而是保证所有产品非零曝光、至少$n-k$个产品取得$\ell$份，同时满足用户EF1与每人恰好$k$项。每个产品在这里被作为生产者侧的一个分配单元。

算法的第一阶段为每个产品提供$\ell$份“可推荐副本”。副本表示可以分给多少位不同用户，不是要把同一个产品重复展示给一个用户。固定一个用户轮次，在轮到$u$时，从尚有副本且$u$尚未获得的产品中，选择$\bm{V}_{u,i}$最高者；分配后减去一份，并从$u$的允许集合中移除该产品。平局按固定规则处理。所有副本耗尽，或当前轮到的用户已经找不到合法产品时，第一阶段终止。

第一阶段终止不表示所有用户都已有$k$项。第二阶段为产品恢复足够的可分配份额，但保留各用户已经获得的产品和去重限制；从下一位尚未获得本轮机会的用户继续循环，直到每人恰好$k$项。更精确地说，若第一阶段在一个不完整轮次中停止，先补完人数较少的这一轮，再完成剩余整轮。沿用轮次位置避免让早选的用户在边界处额外占一次机会。任意跳过无合法对象的用户、任意重启轮次，都不是原算法中可以忽略的改动。

表~\ref{tab:rs-fairrec-example}给出一个三用户、五产品、每人两项的假设例子。此时$\ell=\lfloor6/5\rfloor=1$。第一轮用户1、2、3依次选A、B、C，剩余D、E各一份；第二轮用户1选D、用户2选E，副本用完，用户3仍只有C。第二阶段从用户3继续，允许它取得尚未得到的A，完成分配。三个集合为$\{A,D\}$、$\{B,E\}$和$\{C,A\}$，产品曝光为$(2,1,1,1,1)$。

\begin{table}[htbp]
 \centering
 \small
 \begin{tabular}{@{}lrrrrr@{}}
  \toprule
  用户 & A & B & C & D & E \\
  \midrule
  1 & 9 & 8 & 3 & 2 & 1 \\
  2 & 8 & 9 & 2 & 1 & 3 \\
  3 & 8 & 7 & 9 & 3 & 2 \\
  \bottomrule
 \end{tabular}
 \caption{FairRec教学例子的非负相关性矩阵。列A至E表示五件不同产品；效用取分数之和，用户按1、2、3的固定次序轮流选择。}
 \label{tab:rs-fairrec-example}
\end{table}

用户1所得效用为11，而它对用户3集合的评价为$3+9=12$，所以仍存在普通意义上的嫉妒。移去用户3集合中的A后，剩余效用为3，小于11，因此满足这一对的EF1检查。其他用户对也可用同一矩阵核验。这个例子中所有产品都达到MMS，但它只是某个输入的结果，不能把一般保证中的“至少$n-k$个”加强为“全部”。

生产者侧保证可以从第一阶段的停止条件理解。若全部副本耗尽，每个产品都获得$\ell$份；若某用户$u$无合法产品，它尚未得到的产品一定已经耗尽副本，因此都达到$\ell$份。其已经得到的产品至少有一次曝光，而$u$所得不超过$k$项，所以至少$n-k$个产品达到$\ell$。用户侧EF1则依赖轮流取当前最高效用对象的比较关系：晚取用户可能错过早取用户的第一件对象，但去掉这件对象后，可以按后续轮次配对比较。第二阶段保留轮次衔接与去重，是保持整体分配性质的组成部分。

朴素实现共分配$mk$次，每次扫描$n$个产品，时间为$O(mnk)$；此外要存储或访问相关性矩阵、剩余份额和用户已选集合。它适合讨论一批分配，在大规模实时服务中如何分批、刷新资格、处理动态到达，都需要另行检查。任意用户资格限制、负效用或列表内部互补效应可能破坏原比较条件；只因程序仍能运行，不能继承原保证。

FairRec的输出是用户—产品分配，位置注意仍要交给展示问题处理。一个创作者持有多个产品时，逐产品最低曝光会随作品数量聚合，不能直接称为逐创作者等额公平。MMS也是曝光数量的参照，并不预测曝光能带来多少消费，更不保证创作者继续留在平台。

\subsubsection{Distribution-Free Reliability：集合错误率的风险约束}
\label{sec:rs-distribution-free-reliability}

相关性分配仍依赖模型判断。若一个页面输出十个高分对象，其中有几个其实不符合要求，就需要直接约束最终集合的错误，而不只是约束模型分数。\term{分布无关可靠性方法}（distribution-free reliability）在一个已经训练好的成对排序器外增加集合选择规则，用独立校准请求检验这条规则的平均错误风险。它不预设具体CTR网络，也不要求成对分数本身已经校准成真实概率\scite{rs-e08}

对一个请求，令$C$为含$K\geq2$个对象的候选集，$Y$表示真实次序，$G(Y)\subseteq C$为事先定义的合格集合。例如，在论文的设定中，可以把真实次序前$m$项视为合格；比例和取整规则应在校准前固定。模型给出$\widehat p_{i,j}(\bm{x})$，表示它认为$i$优于$j$的概率。定义质量分数和阈值集合
\begin{equation}
 \begin{aligned}
  q_i(\bm{x})&=\frac{1}{K-1}
           \sum_{j\in C\setminus\{i\}}\widehat p_{i,j}(\bm{x}),\\
  T_\lambda(\bm{x})&=\{i\in C:q_i(\bm{x})\geq\lambda\}.
 \end{aligned}
 \label{eq:rs-risk-set}
\end{equation}
$q_i$汇总成对优势，$\lambda$控制哪些对象被纳入。候选只有一项时，上式分母为零，必须另外定义规则并把它纳入校准，不能默用这个公式。

给定输出集合$S$，\term{错误发现比例}（false discovery proportion, FDP）是其中不合格对象所占的比例：
\begin{equation}
 \operatorname{FDP}(S,Y)
 =\frac{|S\setminus G(Y)|}{\max(1,|S|)},
 \qquad
 R(\lambda)=\mathbb{E}\!
  \left[\operatorname{FDP}(T_\lambda(\bm{X}),Y)\right].
 \label{eq:rs-fdp-risk}
\end{equation}
这里$\bm{X}$表示随机请求信息，$R(\lambda)$又称\term{错误发现率}（false discovery rate, FDR）。分母约定使空集合的FDP为零。期望按请求计算，各请求等权；把所有请求中的错误对象总数除以总推荐数，是另一个聚合量。

例如，假设容易请求的质量分数为$(0.82,0.80,0.20,0.18)$，难请求为$(0.55,0.52,0.49,0.44)$，每组均可来自一个互补的成对概率矩阵。若$\lambda=0.8$，前者输出两项，后者输出空集合；若降到0.5，后者会输出两项。若难请求这两项中只有一项合格，其FDP为$1/2$。模型分数并没有告诉我们在未来请求中这种错误出现多频繁，因此还需要独立数据校准。

设校准集$\mathcal{D}_{\mathrm{cal}}=\{(\bm{x}_r,Y_r)\}_{r=1}^{n}$与模型训练数据分离，并与未来请求独立同分布。给定阈值的经验风险为
\[
 \widehat R_n(\lambda)=\frac1n\sum_{r=1}^{n}
   \operatorname{FDP}(T_\lambda(\bm{x}_r),Y_r).
\]
因为每个损失位于$[0,1]$，Hoeffding不等式可给出固定阈值的一侧上置信界
\begin{equation}
 U_n(\lambda)=
 \min\left\{1,\widehat R_n(\lambda)
       +\sqrt{\frac{\log(1/\delta)}{2n}}\right\}.
 \label{eq:rs-risk-upper-bound}
\end{equation}
$\delta\in(0,1)$控制统计错误的概率，$\alpha\in(0,1)$是允许的风险水平。只有$U_n(\lambda)\leq\alpha$时，才将该阈值通过为风险受控的候选；没有通过不能解释为已经证明风险超标，只表示这份数据不足以确认达标。

阈值由数据选择，不能对许多阈值逐一检验后随意挑最有利的一个，同时沿用单次置信水平。论文采用预定顺序检验（fixed-sequence testing）：校准前固定从高到低的有限阈值序列，依次检验，遇到第一个不能通过的阈值就停止，选择此前最后一个通过者。若首个阈值就失败，返回事先定义、恒为空的集合规则；这个回退应与$\lambda=1$区分，因为模型可能恰好给某对象$q_i=1$。

该程序控制的是
\begin{equation}
 \mathbb{P}_{\mathcal{D}_{\mathrm{cal}}}
 \left(
  \mathbb{E}_{(\bm{X},Y)}
   \left[
    \operatorname{FDP}
     \bigl(T_{\widehat\lambda}(\bm{X}),Y\bigr)
    \mid\mathcal{D}_{\mathrm{cal}}
   \right]>\alpha
 \right)
 \leq\delta.
 \label{eq:rs-reliability-guarantee}
\end{equation}
外层概率反映校准样本的变化，内层期望反映固定校准结果后遇到的新请求。它不保证每次请求的FDP都小于$\alpha$，也不保证每个用户群体的条件风险都达标。

固定顺序为何能避免任意挑选造成的错误累积，可以从第一个真实不达标的规则理解。在校准前已经确定的顺序中，若最后选出了任何不达标规则，就一定通过了这个最先出现的不达标规则；它被错误通过的概率至多为$\delta$。遇到失败后继续寻找后面的偶然低风险点，会破坏这个论证。集合随阈值降低而扩大，但FDP及其期望未必单调，程序的有效性依赖检验顺序和停止方式，而非未经证明的风险单调性。图~\ref{fig:rs-risk-calibration}强调了这一停止条件。

\begin{figure}[htbp]
 \centering
 % 固定顺序检验；失败后停止，不再试探后续阈值。
\begin{tikzpicture}[
 x=1mm,y=1mm,font=\figurefont\small,
 process/.style={draw=BookBlue,fill=BookBlue!5,line width=0.85pt,
                 minimum width=51mm,minimum height=10mm,align=center},
 output/.style={draw=BookTeal,fill=BookTeal!5,line width=0.85pt,
                minimum width=34mm,minimum height=14mm,align=center},
 flow/.style={->,>=latex,draw=BookInk,line width=1.1pt,
              rounded corners=1.2mm}
]
 \path[use as bounding box] (0,0) rectangle (112,85);
 \node[process] (start) at (35,77) {固定模型、规则与阈值顺序};
 \node[process] (sets) at (35,58) {按当前阈值生成校准集合};
 \node[process] (check) at (35,39) {判断：$U_n(\lambda)\leq\alpha$？};
 \node[process] (save) at (35,17) {保存通过的规则};
 \node[output] (stop) at (94,39) {停止\\返回保存规则};
 \draw[flow] (start.south) -- (sets.north);
 \draw[flow] (sets.south) -- (check.north);
 \draw[flow] (check.south) -- node[right,font=\figurefont\footnotesize]
    {通过} (save.north);
 \draw[flow] (check.east) -- node[above,font=\figurefont\footnotesize]
    {失败} (stop.west);
 \draw[flow] (save.west) -- (3,17) -- (3,58) -- (sets.west);
 \node[text=BookMuted,font=\figurefont\footnotesize,rotate=90]
    at (0,38) {仍有阈值：继续};
 \draw[flow] (save.east) -- (94,17) -- (stop.south);
 \node[text=BookMuted,font=\figurefont\footnotesize] at (81,10)
    {序列用尽};
\end{tikzpicture}

 \caption{固定顺序校准集合规则。阈值序列在查看校准误差前确定；每次按完整规则产生集合并计算风险上界。首个检验失败就停止，保留此前通过的规则；若从未通过，则输出预定的空集合回退。}
 \label{fig:rs-risk-calibration}
\end{figure}

用一组假设汇总值检查选择过程。取$n=1000$、$\delta=0.05$，式~\eqref{eq:rs-risk-upper-bound}的误差项约为0.0387。若$\alpha=0.1$，依序考察$\lambda=0.9,0.8,0.7$，经验风险分别为$0.02,0.04,0.08$，上界约为$0.0587,0.0787,0.1187$；前两个通过，第三个失败，因而选择0.8。这里“千个样本”指千个满足独立条件的校准请求，不能把同一请求内高度相关的候选对当成千次独立观测。

可靠性约束还可以包住多样性选择。令$D_\lambda(\bm{x})$从$T_\lambda(\bm{x})$中选至多$M$个对象，按固定的嵌入距离准则提高差异；组合枚举昂贵时，可采用事先确定的贪心近似。应对$D_\lambda$的最终集合重新计算每个校准请求的FDP，并沿同一预定顺序校准。风险保证来自对完整规则的有效检验，不要求多样性求解器达到全局最优。

重新校准不能省略，因为取子集不保证FDP降低。一个集合有九个合格对象和一个不合格对象，原FDP为0.1；若多样性筛选恰好只留下那个不合格对象，FDP变为1。风险通过后再临时改变去重、截断或多样性规则，也会改变被保证的对象。式~\eqref{eq:rs-reliability-guarantee}只约束风险，并不保证所选规则在所有可行规则中实现最大多样性。

计算上，直接汇总全部成对分数需要每个请求$O(K^2)$次成对评价；分数可预先缓存，多个阈值再复用。更复杂的子集选择另计代价。统计上，模型和集合规则在校准前应固定，真实合格标签必须可得；未经处理的选择性点击日志通常不满足这里的标签条件，时间漂移也会破坏同分布假设。平均集合长度、空集合比例、分群错误和实际任务效用应与总体FDR一起报告。一个总是返回空集合的规则风险为零，却没有完成推荐任务。

\section{方案比较与诊断}
\label{sec:rs-decision-diagnostics}

比较方案时，应先固定目标口径、可行范围和证据条件。相同的总收入不能直接比较不同流量规模，相同的曝光份额也不能直接比较不同资格集合。一个方案使用更多候选、更长时间或更多计算预算，其提高可能来自资源增加；若允许这种增加，应把成本同时报告。

对所有目标都越大越好的设定，若方案$a$满足$f_\ell(a)\geq f_\ell(b)$对每个$\ell$成立，且至少一项严格大于，就称$a$支配$b$。没有任何可行方案支配的方案为Pareto最优。这个定义排除了可以不损害其他目标的共同改善，但没有给多方诉求自动确定一个偏好。

本章的A、B、C互不支配：参与度改善都伴随收入下降。再增加D$(0.7,2)$，D会被B$(0.8,3)$支配，应与可行但取舍不同的三个方案区分。若某方案虽有$(0.95,6)$的目标值却超出预算，它不属于可行比较范围。样本中互不支配也只说明没有被当前比较方案支配，不等于已经穷尽全部可行方案，更不等于真实目标在预测误差下也具有同样关系。

严格正权重下，加权目标的全局最优解一定不被支配：若存在支配者，每项贡献不减且至少一项因正权重而增加，会与全局最优矛盾。若某些权重为零，最优解仍可能在被忽略目标上被改进；若只找到局部或近似解，同样需要单独检查。反过来，非凸目标集合中的某些Pareto点无法由固定权重和取得，相关例子见\BookSectionRef{sec:opt-multiobjective-saddle}{多目标与鞍点}。因此，Pareto比较是一种诊断关系，不能当成已经指定的求解算法。

敏感性检查应针对实际规则。A、B、C在$\lambda=0.05$和0.1附近发生选择变化，若候选权重离边界很近，小预测误差就可能改变结果。可以在预先认可的权重、阈值和误差范围内重算可行集合与选择，报告目标变化和资源余量；稳定选择不一定价值最高，却能提示方案对估计误差的脆弱程度。目标阈值没有可行解时，还可逐项放松作诊断，但诊断用的放松不能未经确认就转化为上线承诺。

总体结果也可能掩盖少数群体的损失。假设两个用户群体占比90\%和10\%，基线有效参与率都为0.8。新策略让前者达到0.85、后者降到0.4，则总体为$0.9\times0.85+0.1\times0.4=0.805$，总体小幅上升而小群体明显下降。是否接受这一结果，应由群体目标和约束决定。提供者侧同样需要同时看平均曝光、低机会产品比例、不同规模创作者以及资格变化，不能用总曝光守恒推断分配没有变化。

诊断时可以沿三个不同问题寻找原因。若优化器忠实执行了目标，却反复推荐标题夸张、任务价值低的内容，应检查点击代理是否误设；若目标和约束合理，但预测收入高而实际收入低，应检查概率、价值估计、反馈成熟及分布变化；若预测条件下本可满足预算，实际输出却重复对象或超支，就要检查组合求解和执行状态。三类原因可以同时出现，但分别对应问题定义、学习依据与执行机制，不能一律用增大模型解决。

约束检查还应报告违反程度与证据类型。确定性资格可以逐次检查；期望预算或总体风险则要与其概率保证对应；统计检验未能确认达标，不等于已经证明有害。对整数规划和近似列表构造，还应报告可行解、求解时限及近似误差。跨请求供给响应会进一步改变目标和可行范围，由第~\ref{chap:rs-long-term}~章扩展；离线与在线证据怎样识别实际效果，由第~\ref{chap:rs-evaluation}~章建立。

\section{本章小结}
\label{sec:rs-constraints-summary}

一次推荐是否值得执行，需要同时回答多方希望得到什么，以及资格、资源和底线允许选择什么。目标值只有在主体、单位、聚合方式与时间范围一致时才可比较。加权规定交换率，优先级规定比较次序，目标约束规定可接受范围；它们可以组合，但训练权重、业务底线和实际执行保证不能互相替代。

公平与可靠性把这种区别具体化了。累计注意分配以相关性比例和位置模型为依据，FairRec在可加效用和集合分配条件下保证EF1与部分产品的MMS份额，集合可靠性则用独立校准数据控制未来请求的平均错误比例。三者约束不同对象，也需要不同证据。得到一个可行方案后，还要检查支配关系、分群后果、敏感性和计算代价。

本章已经能够利用累计状态限制当前行动，却没有完整回答当前行动怎样改变未来。曝光可能增加疲劳，也可能补充信息；提供者还会据机会改变供给。下一章把这些变化放进状态转移与累计价值，使本章的目标和资源约定成为连续决策的一部分。
